\section{The Philosophy of Well-Being and the Theory of Justice}
\label{sec:philosophy}

\noindent
The paper's conception of a just eudaimonia stands in two philosophical lineages that a foundational treatment must make explicit: the philosophy of well-being, which asks what a good life consists in, and the theory of justice, which asks what is owed among those who share a life. This section locates the conception within each, not to derive it from them, since it was developed on the independent grounds of the prior papers, but to situate it, to mark its agreements and departures, and to make available to the later argument the distinctions these literatures have drawn.

\subsection{The philosophy of well-being}
\label{sec:phil-wellbeing}

The philosophy of well-being distinguishes three families of theory of what makes a life go well, a division owed in its canonical form to Parfit \citep{parfit1984}. Hedonistic theories hold that well-being consists in the balance of pleasure over pain, of positive over negative experience; the good life is the happiest life, and the tradition descends through the classical utilitarians and Sidgwick \citep{sidgwick1907}. Desire-fulfilment theories hold that well-being consists in the satisfaction of a person's desires, whether or not the satisfaction is experienced as pleasant; the good life is the life in which what one wants is achieved, and the theory has the advantage of grounding a person's good in that person's own valuings rather than in an external standard. Objective-list theories hold that certain things are good for a person whether or not they are desired or found pleasant, that there is a plurality of goods, knowledge, friendship, achievement, virtue, and the like, whose possession makes a life better independently of the person's attitude toward them; the good life is the life furnished with the items on the list.

The conception of this series stands in a determinate relation to these three. It is not hedonistic: the paper has argued at length that the felt satisfaction of a life is a structured proxy that a forged flourishing can satisfy (\xref{sec:ref-blind}), and that a relation may produce a high balance of positive affect while being, in its distribution of return, unjust and therefore no eudaimonia. It is not a simple desire-fulfilment theory: the paper has held that the value at issue is not exhausted by what the parties happen to want, since a party may, through habituation to a diminished position, cease to want what the foreclosure of their unfolding has put beyond reach, so that the satisfaction of their now-diminished desires is compatible with the foreclosure of their flourishing. The conception is closest to an objective-list theory, in that it holds the unfolding of each party's own dynamics, and the justice of the cycle that sustains it, to be good for the parties whether or not they are felt or desired; but it differs from the standard objective-list theory in two respects. It does not offer a list of goods but a structure, the open spiral of a returning cycle, within which the goods of a particular shared life are generated rather than enumerated in advance; and it grounds the objectivity of this structure not in an external standard imposed on the parties but in the dynamics of generation and return that the prior papers derived, so that what is good for the parties is what their own generative activity, justly sustained, produces. The conception is, in the vocabulary of the philosophy of well-being, an objective-list theory whose single item is a structure rather than a content, and whose objectivity is that of a generative process rather than of an imposed standard.

\subsection{The theory of justice}
\label{sec:phil-justice}

The condition of justice that the series holds to be constitutive of eudaimonia stands in relation to the principal contemporary theories of justice, and the relation illuminates both what the paper's justice condition is and what it adds. The dominant modern theory of distributive justice, owed to Rawls, concerns the fair distribution of primary goods, the goods any rational person is presumed to want, among them not only income and opportunity but the social bases of self-respect, and it assesses a social arrangement by the principles that would govern this distribution were they chosen from behind a veil of ignorance \citep{rawls1971}. The paper's justice condition is distributive in form, since it concerns the distribution of returned value among the parties of a relation, and the skew of the return distribution (\xref{sec:ev-skew}) is a measure of distributive injustice within the relation; the inclusion, among Rawls's primary goods, of the social bases of self-respect is of particular pertinence, since the party consigned to the low-return tail is deprived not only of returned value but of the self-respect that being returned to sustains.

A second tradition, relational egalitarianism, holds that the concern of justice is not in the first instance the distribution of goods but the character of the relations in which people stand, that justice obtains when people relate as equals rather than when goods are distributed by a pattern, and that the point of equality is the abolition of relations of domination, exploitation, and subordination \citep{anderson1999}. The paper's justice condition is, in this respect, relational as much as distributive: the injustice it identifies is the reduction of one party to the sustaining fuel of a relation, a relation of exploitation and subordination within the shared life, and the skew of return is the signature of such a relation and not merely of an unequal distribution. The two construals, the distributive and the relational, are for the paper two aspects of one condition, the distribution of return being the measurable trace of the relational standing of the parties.

A third tradition supplies the connection to recognition. The debate between redistribution and recognition, joined by Fraser and Honneth, concerns whether injustice is fundamentally a matter of maldistribution, to be remedied by redistributing goods, or of misrecognition, the failure to accord a person or group due standing, to be remedied by recognition \citep{fraserhonneth2003}. The paper's account integrates the two: the return of value to its generator is at once a redistribution, the circulation of value back to the party who produced it, and a recognition, the acknowledgement of that party as a generator whose contribution is owed a return rather than a source from which value may be taken without acknowledgement. The skew of return measures a maldistribution; the foreclosure of unfolding (\xref{sec:ev-dynamics}) measures a misrecognition, the failure to accord a party the standing of one whose own development matters; and the paper holds these to be two faces of a single injustice, as the redistribution-recognition debate, in its more integrative resolutions, has come to hold.

A fourth tradition, the ethics of care, is the one nearest the paper's subject matter and the one whose insight the paper most directly formalises. The care tradition, arising from Gilligan's identification of a moral voice oriented to responsibility and relationship rather than to abstract right \citep{gilligan1982}, and developed into a political and philosophical ethics by Held and others \citep{held2006}, holds that the care through which dependents are sustained is morally central and systematically undervalued, and that a just account of social life must reckon the labour of care as labour rather than effacing it as natural or freely given. Kittay's analysis of dependency and the labour it requires is of particular pertinence, since it identifies exactly the sustaining labour whose under-return the paper's skew measure is constructed to detect \citep{kittay1999}. The paper may be read as supplying, for the central case of the intimate relation, a formalisation of the care tradition's core claim: that the labour of sustaining a shared life is labour, that it generates value, and that its systematic under-return is an injustice with a measurable signature. The ethics of care supplies the moral insight; the skew of the return distribution offers a structured proxy for its violation.

\subsection{The location of the paper's conception}
\label{sec:phil-location}

The paper's conception of a just eudaimonia is, in sum, an objective-list theory of well-being whose single item is the structure of a justly returning cycle, joined to a theory of justice that is distributive in its measure, relational in its content, integrative of redistribution and recognition, and formalising of the ethics of care. The conception was not assembled from these positions but developed on independent grounds, and its relation to them is one of convergence rather than derivation: the prior papers' account of generative justice, the return to the creator, and the unfolding of each arrived, by its own route, at a conception that the philosophy of well-being and the theory of justice illuminate from their several directions. The convergence is itself an argument of a kind, since a conception reached independently and found to gather the insights of several established traditions has a claim to have articulated something the traditions severally approach; but the paper rests its conception on the prior derivation and offers the convergence as corroboration and not as ground.
